\documentclass[10pt,a4paper]{amsart}
\usepackage[margin=19mm]{geometry}
\usepackage{amsmath,amssymb,mathtools}
\usepackage[scr=rsfs,cal=euler]{mathalfa}
\usepackage{xfrac}
\usepackage[unicode,hidelinks,bookmarks=false]{hyperref}
\newcommand{\ie}{i.e.\ }
\newcommand{\cf}{cf.\ }
\newcommand{\resp}{respectively\ }
\newcommand{\Subsection}{\subsection}
\providecommand{\nobreakdash}{\nobreak\hbox{-}}
\setlength{\emergencystretch}{2em}
\multlinegap0pt
\usepackage{amssymb}
\usepackage{enumerate}
\usepackage{xcolor}
\newtheorem{cor}{Corollary}
\newtheorem{conj}{Conjecture}
\newtheorem{lem}{Lemma}
\usepackage{cleveref}
\crefname{cor}{Corollary}{Corollarys}
\crefname{conj}{Conjecture}{Conjectures}
\crefname{lem}{Lemma}{Lemmas}
\newcommand{\Dk}{\mathcal D_k}
\newcommand{\Dtk}{\widetilde{\mathcal D}_k}
\DeclareMathOperator{\Fil}{Fil}
\newcommand{\ZAfw}[1]{{\mathcal Z_{\mathcal A,#1}^{\mathrm f}}}
\newcommand{\za}[1]{\zeta_{\mathcal A}(#1)}
\newcommand{\zaf}[1]{\zeta_{\mathcal A}^{\mathrm f}(#1)}
\newcommand{\zzf}[1]{\mathfrak z^{\mathrm f}(#1)}
\title{Formal finite multiple zeta values}
\author{Henrik Bachmann, Risan}
\date{Source: arXiv:2608.15480v1 (CC BY 4.0)}
\begin{document}
\maketitle

\noindent\emph{Source and licence.} Formal finite multiple zeta values. Version arXiv:2608.15480v1 (CC BY 4.0), distributed by arXiv under the Creative Commons Attribution 4.0 International licence, http://creativecommons.org/licenses/by/4.0/, as recorded in the arXiv licence metadata for this exact version and shown on the abstract page https://arxiv.org/abs/2608.15480v1. Full licence notice: \texttt{licenses/arxiv-2608.15480-CC-BY-4.0.txt}.\par\smallskip

\noindent\textit{Editorial scope.} Counted unit: the article's lem lem:direct-beta-certificate (source0.tex lines 1643--1650). The article's complete original proof of it is delivered with it (source0.tex lines 1652--1781, one \texttt{proof} environment of 130 lines). No mathematics is written, repaired or re-derived: the statement and the proof are the article's own lines, in the article's own notation, and no retained line is substituted or re-flowed.  Retained uncounted as the source-local context the proof needs: lines 274--281; lines 1100--1115; lines 1150--1154; lines 1437--1442; lines 1453--1465; lines 1474--1480; lines 1620--1628.  Nothing was omitted from any retained span, so the package carries no omission record.  No retained line contains a citation, so the wrapper carries no bibliography and no bibliography entry.  No retained line contains a figure environment, an image-inclusion command or drawing code, so no image is delivered and the package declares no asset dependency.  Editorial formatting only, in addition: an AMS \texttt{amsart} preamble that restates the article's own declarations and macros used by the retained lines, namely the environments \texttt{cor}, \texttt{conj}, \texttt{lem} on separate counters, together with the standard packages those lines need.  The retained environments print in this document as Corollary 1, Conjecture 1, Lemma 1 in order of appearance, whereas the article prints the same items with the numbering of its own full document.  The complete native source of the article as distributed in the arXiv e-print for this exact version is bundled unchanged as \texttt{sources/207751/source0.tex}.
\begin{align}\label{eq:H-source-element}
 H_{r,s}=\frac12\Bigg(&
 (rs-r+1)Z_{r,s}+(rs+s-1)Z_{s,r}
 +(s+1)Z_{s+1,r-1}-(r+1)Z_{r+1,s-1}\notag\\
 &+\sum_{\substack{3\leq\ell\leq k-3\\\ell\ \mathrm{odd}}}
 \left(\binom{k-\ell}{s}-\binom{k-\ell}{r}\right)
 P_{\ell,k-\ell}\Bigg)\in\Dtk.
\end{align}

\begin{cor}
\label{cor:explicit-true-single-product}
Let $r,s\geq3$ be odd and put $k=r+s$. Then
\begin{align}\label{eq:four-one-product}
 \zaf{r-1,1,s-1,1}
 =\frac12\zaf{r-1,1}\zaf{s-1,1}
  +\frac12\bigl(\zaf{1,r-1,s}-\zaf{1,s-1,r}\bigr).
\end{align}
Consequently,
\begin{align}\label{eq:four-one-symmetry}
 \zaf{r-1,1,s-1,1}+\zaf{s-1,1,r-1,1}
 &=\zaf{r-1,1}\zaf{s-1,1},\\
 \zaf{r-1,1,s-1,1}-\zaf{s-1,1,r-1,1}
 &=\zaf{1,r-1,s}-\zaf{1,s-1,r}.\notag
\end{align}
\end{cor}

\begin{align}\label{eq:formal-double-shuffle}
 P_{r,s}
 &=Z_{r,s}+Z_{s,r}+Z_k,\\
 &=\sum_{a+b=k}\left[\binom{a-1}{r-1}+\binom{a-1}{s-1}\right]Z_{a,b}.\notag
\end{align}

\begin{align}\label{eq:direct-beta-rules}
 \beta_k^{\mathrm f}(H_{r,s})
 &=\zaf{r-1,1,s-1,1},\notag\\
 \beta_k(H_{r,s})
 &=\za{r-1,1,s-1,1}
\end{align}

\begin{conj}\label{conj:beta-identities}
For every even $k\geq4$ and all $a,b,c\geq1$ with $a+b+c=k$, one has
in $\ZAfw{k}$
\begin{align}\label{eq:formal-Hoffman-depth-three}
 \zaf{a,b,c}
={}&(-1)^{b-1}
 \sum_{\substack{n+m=k-1\\n,m\geq1}}
 \binom na\binom mc\,\zaf{n,1,m}\notag\\
&+(-1)^{a+c}\sum_{i=1}^{b-1}
 \binom{a+i}{a}\binom{b+c-i}{c}
 \zzf{a+i}\zzf{b+c-i}.
\end{align}
\end{conj}

\begin{align}\label{eq:formal-harmonic}
0=
 \sum_{\substack{i+j=r,\ a+b=s\\i,j,a,b\geq1\\i+a\ \mathrm{odd}}}
 \zaf{i,a}\zaf{j,b}
 +2(-1)^r\sum_{\substack{u+v=k-1\\u,v\geq1}}
 \left(\binom vr+\binom vs\right)\zaf{u,1,v}.
\end{align}

\begin{lem}\label{lem:restricted-sum}
For every even $k\geq4$,
\begin{align}\label{eq:restricted-sum}
 \sum_{a+b=k-1}\zaf{1,a,b}
 =\sum_{a+b=k-1}\zaf{a,1,b}
 =\sum_{a+b=k-1}\zaf{a,b,1}=0,
\end{align}
where the summation variables are positive.
\end{lem}

\begin{lem}\label{lem:direct-beta-certificate}
Assume that \cref{conj:beta-identities} holds in weight $k$. Then the
first rule in \eqref{eq:direct-beta-rules} defines a linear map
\[
 \beta_k^{\mathrm f}:\Dtk\longrightarrow\Fil_4\ZAfw{k}.
\]
Its image is $\Fil_3\ZAfw{k}$.
\end{lem}

\begin{proof}
We verify compatibility by lifting the prescribed values to the standard
source generators. Put
\[
 \mu_k=\frac12
 \sum_{\substack{1\leq n\leq k-1\\n\ \mathrm{odd}}}
 \zzf n\zzf{k-n}.
\]
For $1\leq r\leq k-1$, define
\begin{align}\label{eq:direct-beta-certificate}
 y_r=(-1)^r\sum_{n=1}^{k-2}
 \left(\binom nr-\frac12\binom{k-1}r\right)
 \left(\zaf{n,1,k-n-1}-\zzf n\zzf{k-n}\right)
 -\delta_{r,1}\mu_k.
\end{align}

Assign $Z_k\mapsto0$ and $Z_{r,k-r}\mapsto y_r$ for
$1\leq r\leq k-1$, and abbreviate
$D_n=\zaf{n,1,k-n-1}-\zzf n\zzf{k-n}$.
We first evaluate the shuffle side
$\sum_{a=1}^{k-1}\bigl[\binom{a-1}{r-1}+\binom{a-1}{s-1}\bigr]y_a$ of
\eqref{eq:formal-double-shuffle}, where $s=k-r$. The inversion
identity
\[
 \sum_{a\geq1}(-1)^a\binom{a-1}{r-1}\binom na=
 \begin{cases}
  (-1)^r,&n\geq r,\\
  0,&n<r,
 \end{cases}
\]
applied in \eqref{eq:direct-beta-certificate} gives
\[
 \sum_{a=1}^{k-1}\binom{a-1}{r-1}y_a
 =(-1)^r\sum_{n=r}^{k-2}D_n
 -\frac{(-1)^r}2\sum_{n=1}^{k-2}D_n
 -\delta_{r,1}\mu_k.
\]
By \cref{lem:restricted-sum} and reversal,
\begin{align}\label{eq:boundary-sums}
 \sum_{n=1}^{k-2}D_n=-2\mu_k,
 \qquad
 \sum_{n=r}^{k-2}D_n+\sum_{n=s}^{k-2}D_n=-2\mu_k-\zzf r\zzf s,
\end{align}
so the shuffle side equals
$-(-1)^r\zzf r\zzf s-(\delta_{r,1}+\delta_{s,1})\mu_k$, that is,
$\zzf r\zzf{k-r}$ for $2\leq r\leq k-2$ (both sides vanish for even
$r$) and $-\mu_k$ for $\{r,k-r\}=\{1,k-1\}$. On the other hand, a
Vandermonde convolution
$\sum_i\binom ui\binom v{r-i}=\binom kr$,
followed by \cref{lem:restricted-sum}, rewrites
\eqref{eq:formal-harmonic} as
\begin{align}\label{eq:harmonic-binomial-direct}
 \sum_{n=1}^{k-2}
 \left(\binom nr+\binom ns-\frac12\binom kr\right)D_n
 =-\zzf r\zzf s.
\end{align}
Since
$\binom{k-1}r+\binom{k-1}{s}=\binom kr$, comparing with
\eqref{eq:direct-beta-certificate} gives
$y_r+y_{k-r}=\zzf r\zzf{k-r}$ for
$2\leq r\leq k-2$. At the boundary,
\cref{lem:restricted-sum} and reversal give
\[
 \sum_{n=1}^{k-2}\left(n-\frac{k-1}{2}\right)D_n=-\mu_k,
\]
which together with the first sum in \eqref{eq:boundary-sums},
substituted in \eqref{eq:direct-beta-certificate}, gives
\[
 y_1=0,\qquad y_{k-1}=-\mu_k.
\]
Thus the stuffle and shuffle values agree in every case. Sending each
product symbol to this common value defines a realization of $\Dk$.
By construction $Z_k\mapsto0$, and the boundary calculation gives
$Z_{1,k-1}\mapsto y_1=0$. Each even product symbol $P_{r,k-r}$ maps to
$\zzf r\zzf{k-r}=0$, so the realization descends to $\Dtk$.

Moreover,
\[
 P_{r,k-r}\longmapsto\zzf r\zzf{k-r}
 \qquad(2\leq r\leq k-2).
\]
By
\eqref{eq:four-one-product},
\begin{align}\label{eq:four-one-bridge}
 2\zaf{r-1,1,s-1,1}={}&rs\zzf r\zzf s\notag\\
 &+\zaf{1,r-1,s}-\zaf{1,s-1,r}.
\end{align}
Substituting \eqref{eq:direct-beta-certificate} in
\eqref{eq:H-source-element}, using
$y_r+y_s=\zzf r\zzf s$, and applying
\eqref{eq:formal-Hoffman-depth-three} with
$(a,b,c)=(1,r-1,s)$ and $(1,s-1,r)$ gives
\[
 H_{r,s}\longmapsto\frac12\left(
 rs\zzf r\zzf s+\zaf{1,r-1,s}-\zaf{1,s-1,r}
 \right).
\]
By \eqref{eq:four-one-bridge}, this is
$\zaf{r-1,1,s-1,1}$. Thus the descended map is the map prescribed
in \eqref{eq:direct-beta-rules}.

It remains to determine its image. Add $\mu_k$ back to the first
coordinate in \eqref{eq:direct-beta-certificate}, and denote the
resulting coordinates by $x_r$. Since $x_{k-1}=-\mu_k$, the spans of
the $x_r$ and the $y_r$ agree. For $1\leq r\leq k-2$, put
\[
 U_r=(-1)^rx_r+\binom{k-1}{r}x_{k-1}.
\]
Then
\[
 U_r=\sum_{n=1}^{k-2}\binom nr
 \left(\zaf{n,1,k-n-1}-\zzf n\zzf{k-n}\right),
\]
and binomial inversion gives
\[
 \zaf{n,1,k-n-1}-\zzf n\zzf{k-n}
 =\sum_{r=n}^{k-2}(-1)^{r-n}\binom rnU_r.
\]
Thus every difference on the left belongs to the image. For
$2\leq n\leq k-2$, \eqref{eq:harmonic-binomial-direct} gives
$\zzf n\zzf{k-n}=x_n+x_{k-n}$. Hence
$\zaf{n,1,k-n-1}$ belongs to the image. This is also true for $n=1$,
since $\zzf 1=0$. Finally,
\eqref{eq:formal-Hoffman-depth-three} expresses every depth-three value
in terms of these values and the products $\zzf n\zzf{k-n}$. The image
therefore contains $\Fil_3\ZAfw{k}$. Conversely, the image is spanned
by the values $\zaf{r-1,1,s-1,1}$, which belong to $\Fil_3\ZAfw{k}$ by
\cref{cor:explicit-true-single-product}. The image is therefore
$\Fil_3\ZAfw{k}$.
\end{proof}

\end{document}
